\section{Function Encoding}

\red{We investigate how Boolean functions can be encoded in quantum states using \computationCircuits{}.}

Using the \computationCircuit{} we can prepare two function encoding schemes, depending on the start state.


\begin{remark}[Basis and Amplitude Encodings in Quantum Computation \cite{schuld_machine_2021} and \tnreason{}]
    % Compare with Quantum Computing
    Basis Encoding in Quantum Computation refers to the representation of classical $n$ bit strings by $n$ qubit basis states, and is called one-hot encoding in \tnreason{}.
    The Basis Encoding scheme in \tnreason{} goes beyond this scheme and also encodes subsets by sums of one-hot encodings to the members of the set.
    In this way, relations and functions are represented by boolean tensors and contraction of them is refered as \BasisCalculus{}.

    Amplitude Encoding in Quantum Computation refers to the storage of complex numbers in the amplitudes of quantum states.
    The pendant in \tnreason{} is the Coordinate Encoding scheme, where real numbers are stored in the coordinates of real-valued tensors.
    Compared to Amplitude Encoding, Coordinate Encoding does not have the normalization constraint of quantum states.
    The Amplitude Encoding of the square root of a probability distribution is sometimes called q-sample.
\end{remark}


\subsection{Basis encoding}

\begin{definition}
    The basis encoding of a boolean function $\exformula$ is the state
    \begin{align*}
        \basqstateofat{\exformula}{\shortcatvariables,\headvariable}
        = \sqrt{\frac{1}{2^\atomorder}} \sum_{\shortcatindices\in\atomstates} \onehotmapofat{\shortcatindices}{\shortcatvariables}
        \otimes \onehotmapofat{\exformula(\shortcatindices)}{\headvariable} \, .
    \end{align*}
\end{definition}

\begin{lemma}
    When applying the \computationCircuit{} on the initial state
    \begin{align*}
        \sqrt{\frac{1}{2^\atomorder}}\onesat{\shortcatvariables} \otimes \fbasisat{\headvariable} % insymbol?
    \end{align*}
    we get the basis encoding state (see \figref{fig:basSignQstateConstruction}a).
\end{lemma}
\begin{proof}
    Since the Pauli-X turns $\fbasis$ into $\tbasis$, and is applied exactly when $f(\shortcatindices)=1$.
\end{proof}


\begin{corollary}
    Using \lemref{lem:bencComCircuit} we have
    \begin{align*}
        \basqstateofat{\exformula}{\shortcatvariables,\headvariable}
        = \sqrt{\frac{1}{2^{\catorder}}} \bencodingofat{\exformula}{\shortcatvariables,\headvariable}
    \end{align*}
\end{corollary}

\begin{figure}[t]
    \begin{center}
        \input{tikz/function_encoding/bas_sign_qstate_construction}
    \end{center}
    \caption{Construction of the basis encoding state (a) and the sign encoding state (b) using the \computationCircuit{} of a function.}\label{fig:basSignQstateConstruction}
\end{figure}

The initial state can be prepared by applying Hadamard gates on all distributed variable qubits and preparing the head variable qubit in the ground state.
The overlap of two basis encoding states (which share the same head variable) is proportional to the contraction of the corresponding formulas.

\begin{lemma}
    \label{lem:basQstateOverlap}
    Let $\formula,\secformula$ be two Boolean formulas on the states of $\shortcatvariables$.
    Then we have
    \begin{align*}
        \contraction{\basqstateofat{\formula}{\shortcatvariables,\headvariable},\basqstateofat{\secexformula}{\shortcatvariables,\headvariable}}
        = \frac{1}{2^{\catorder}} \contraction{\formulawith,\secformulawith} \, .
    \end{align*}
\end{lemma}
\begin{proof}
    We have
    \begin{align*}
        \contraction{\basqstateofat{\formula}{\shortcatvariables,\headvariable},\basqstateofat{\secexformula}{\shortcatvariables,\headvariable}}
        &= \frac{1}{2^{\catorder}}
        \sum_{\shortcatindices\in\shortatomstates} \contraction{\onehotmapofat{\formulaat{\indexedshortcatvariables}}{\headvariable},\onehotmapofat{\secformulaat{\indexedshortcatvariables}}{\headvariable}} \\
        &= \frac{1}{2^{\catorder}} \contraction{\formulawith,\secformulawith} \, .
    \end{align*}
\end{proof}

\subsection{Sign encoding}

\begin{definition}
    The sign encoding of a boolean function $\exformula$ is the state
    \begin{align*}
        \signqstateofat{\exformula}{\shortcatvariables,\headvariable}
        = \left(\sqrt{\frac{1}{2^\atomorder}} \sum_{\shortcatindices\in\atomstates} \left(-1\right)^{\exformula(\shortcatindices)} \cdot \onehotmapofat{\shortcatindices}{\shortcatvariables}\right)
        \otimes \sqrt{\frac{1}{2}} \left(\fbasisat{\headvariable} - \tbasisat{\headvariable}\right) \, .
    \end{align*}
\end{definition}

Sign encodings can be prepared by the phase kickback trick, as we show next.

\begin{lemma}[Phase kickback trick]
    Applying the \computationCircuit{} $\qcbencodingofat{\formula}{\cheadvariable{\insymbol},\cheadvariable{\outsymbol},\shortcatvariables}$ on a state
    \begin{align*}
        \qstateat{\shortcatvariables} \otimes \hminusbasat{\cavariable{\insymbol}}
    \end{align*}
    we get
    \begin{align*}
        \sum_{\shortcatindicesin} (-1)^{\formula(\shortcatindices)} \qstateat{\indexedshortcatvariables}\cdot \onehotmapofat{\shortcatindices}{\shortcatvariables}
        \otimes \hminusbasat{\cavariable{\outsymbol}} \, .
    \end{align*}
\end{lemma}
\begin{proof}
    We get a disentangled ancilla qubit, because $\hminusbasat{\cavariable{\insymbol}}$ is an eigenvector of the identity and the Pauli-X.
    For the identity, which is applied when $\formula(\shortcatindices)=0$, the eigenvalue is $1$.
    For the Pauli-X, which is applied when $\formula(\shortcatindices)=0$, the eigenvalue is $(-1)$.
\end{proof}

We get the sign encoding with the phase kickback trick, when we initialize with the q-sample of the uniform distribution.

\begin{corollary}% Generalize above: Does not have to be uniform init state !
    When applying the \computationCircuit{} on the initial state
    \begin{align*}
        \sqrt{\frac{1}{2^\atomorder}}\onesat{\shortcatvariables}
        \otimes \sqrt{\frac{1}{2}} \left(\fbasisat{\headvariable} - \tbasisat{\headvariable}\right)
    \end{align*}
    we get the sign encoding state (see \figref{fig:basSignQstateConstruction}b).
    Note that the initial state is a Hadamard gate acting on the state $\tbasisat{\headvariable}$.
\end{corollary}

The sign encoding is also called the phase encoding.
% Other phase encodings
One can ask, whether we can use other ancilla states and get different phase kickbacks while keeping the ancilla disentangled.
However, since we construct here \computationCircuit{} we are restricted to the eigenvectors of the Pauli-X gate, which is the Hadamard basis
\begin{align*}
    \hminusbasat{\headvariable}
    &= \sqrt{\frac{1}{2}} \left(\fbasisat{\headvariable} - \tbasisat{\headvariable}\right) \quad,\quad
    \hplusbasat{\headvariable}
    &=\sqrt{\frac{1}{2}} \left(\fbasisat{\headvariable} + \tbasisat{\headvariable}\right)
\end{align*}
with eigenvalues $(-1)$ and $0$.


% Preparation of the initial state
The initial state can be prepared by applying Hadamard gates on all distributed variable qubits and preparing the head variable qubit by applying a Hadamard gate on the first basis state, i.e.
\begin{align*}
    \frac{1}{\sqrt{2}} \left( \tbasisat{\headvariableof{\outsymbol}} - \fbasisat{\headvariableof{\outsymbol}} \right)
    = \hgateat{\headvariableof{\outsymbol},\headvariableof{\insymbol}} \fbasisat{\headvariableof{\insymbol}} \, ,
\end{align*}

\begin{lemma}
    \label{lem:signQstateOverlap}
    Let $\formula,\secformula$ be two Boolean formulas on the states of $\shortcatvariables$.
    Then we have
    \begin{align*}
        \contraction{\signqstateofat{\formula}{\shortcatvariables,\headvariable},\signqstateofat{\secexformula}{\shortcatvariables,\headvariable}}
        = \frac{1}{2^{\catorder}} \left(\contraction{\formulawith,\secformulawith}-\contraction{\formulawith,\lnot\secformulawith}\right) \, .
    \end{align*}
\end{lemma}
\begin{proof}
    We have
    \begin{align*}
        \contraction{\signqstateofat{\formula}{\shortcatvariables,\headvariable},\signqstateofat{\secexformula}{\shortcatvariables,\headvariable}}
        &= \frac{1}{2^{\catorder}}
        \sum_{\shortcatindices\in\shortatomstates} (-1)^{\formulaat{\indexedshortcatvariables}+\secformulaat{\indexedshortcatvariables}}  \\
        &= \frac{1}{2^{\catorder}}
        \sum_{\shortcatindices\in\shortatomstates \wcols \formulaat{\indexedshortcatvariables}=\secformulaat{\indexedshortcatvariables}} 1
        + \sum_{\shortcatindices\in\shortatomstates \wcols \formulaat{\indexedshortcatvariables}\neq\secformulaat{\indexedshortcatvariables}} (-1) \\
        &= \frac{1}{2^{\catorder}} \left(\contraction{\formulawith,\secformulawith}-\contraction{\formulawith,\lnot\secformulawith}\right) \, .
    \end{align*}
\end{proof}

Limitations:
\begin{itemize}
    \item The prepared state $\signqstateof{\exformula}$ is distinguished from the state $\signqstateof{\lnot\exformula}$ only by a non-observable phase factor
    \begin{align*}
        \signqstateof{\exformula}
        = (-1) \cdot \signqstateof{\lnot\exformula}
    \end{align*}
    \item With this construction the statistic qubit is disentangled with the distributed qubits.
    Further manipulation of the statistic qubits alone (as we do with \activationCircuits{}) will not retrieve any information about the encoded function.
    \item When doing the sign encoding procedure for multiple formulas, we get a disentangled state of the head variables with the distributed variables encoded in the sign encoding state of $\bigoplus_{\exformulain}\exformula$ (where by $\bigoplus$ we denote the mod2 sum of boolean functions).
\end{itemize}


% Multiple formulas

\subsection{Conclusion}

We have seen that preparing some classically hard to compute contraction can be done efficiently by a quantum circuit.
We showed in particular that contracted basis encodings can be prepared by circuits, which are classically hard to contract in the presence of undirected loops in the decomposition hypergraph.
However, when assessing the prepared contraction state by computational basis measurements, we effectively draw samples from the corresponding distribution, which can also be achieved classically.
A quantum advantage could arise from different measurement schemes studied in the quantum state tomography community (see \cite{gross_quantum_2010,cramer_efficient_2010,roth_semi-device-dependent_2023}).
To this end, future research could show, how hard to achieve expectation queries can efficiently be prepared by customized measurement schemes.